\chapter{Time: the transient heat equation with BDF1 and BDF2}
\label{ch:C1_heat_bdf}
\noindent\texttt{tutorials/C\_time/C1\_heat\_bdf.py}\\[0.75em]
\begin{outcome}
You can discretize du/dt with backward-difference
formulas, run a time loop with history rotation, and — the discipline of
this course — MEASURE the temporal orders: 1 for BDF1, 2 for BDF2. You
also learn the two classic traps of time-convergence studies: (i) the
first step of BDF2 has no history (bootstrap with BDF1 — one step at
local O(dt\textasciicircum{}2) does not spoil global order 2), and (ii) spatial error
pollutes dt-ladders unless you compare against a fine-dt reference ON THE
SAME MESH.
\end{outcome}

u\_t - lap(u) = f,  u = g on the boundary, u(0) = u0. The BDF
discretization turns each step into a Poisson-like solve:
    (b0/dt) M u\textasciicircum{}\{n+1\} + K u\textasciicircum{}\{n+1\} = F(t\textasciicircum{}\{n+1\}) - M (b1 u\textasciicircum{}n + b2 u\textasciicircum{}\{n-1\})/dt
with the production tables \{1,-1,0\} (BDF1) and \{1.5,-2,0.5\} (BDF2) from
diffsim.solvers.timestepping. MMS in time: u* = cos(2 pi t) sin(pi x)
sin(pi y), f = u*\_t - lap(u*).

\begin{expected}
\begin{verbatim}
BDF1: rates -> ~1.0        BDF2: rates -> ~2.0
    and BDF2's error at dt = T/32 beats BDF1's at dt = T/128.
\end{verbatim}
\end{expected}

\section*{The script}
\begin{lstlisting}
import numpy as np
import scipy.sparse as sp
from scipy.sparse.linalg import splu
import warp as wp

from diffsim.octree.build import build_uniform
from diffsim.mesh.nodes import build_mesh
from diffsim.mesh.constraints import build_constraints
from diffsim.mesh.basis import basis_tables
from diffsim.assembly.operators import DeviceMesh, assemble_csr
from diffsim.physics.poisson import gauss_points, make_load_kernel
from diffsim.solvers.timestepping import bdf_coeffs, bdf_order_now, History

DEVICE = "cuda:0"
OM = 2 * np.pi
u_ex = lambda x, t: np.cos(OM * t) * np.sin(np.pi * x[:, 0]) \
    * np.sin(np.pi * x[:, 1])
f_ex = lambda x, t: (-OM * np.sin(OM * t)
                     + 2 * np.pi ** 2 * np.cos(OM * t)) \
    * np.sin(np.pi * x[:, 0]) * np.sin(np.pi * x[:, 1])


class HeatSolver:
    def __init__(self, level):
        tree = build_uniform(level, dim=2)
        self.mesh = build_mesh(tree, p=1)
        self.cons = build_constraints(self.mesh)
        self.dm = DeviceMesh.from_mesh(self.mesh, self.cons,
                                       basis_tables(1, dim=2), DEVICE)
        self.K = assemble_csr(self.dm)
        self.T = self.cons.T.tocsr()
        self.M = self._mass()
        self.xq = gauss_points(self.mesh, self.dm.tables_by_p)
        self.coords = self.mesh.node_coords[self.cons.free_nodes]
        self.bdry = np.where(
            self.mesh.boundary_nodes[self.cons.free_nodes])[0]

    def _mass(self):
        dm, mesh = self.dm, self.mesh
        rows, cols, vals = [], [], []
        for pv in dm.bins:
            tb = dm.tables_by_p[pv]
            h = mesh.tree.h()[mesh.bins[pv]]
            jac = (h / 2.0) ** dm.dim
            Me = np.einsum("qa,qb,q->ab", tb.N, tb.N, tb.w)
            conn = mesh.conn_of[pv].astype(np.int64)
            nbf = conn.shape[1]
            Mee = Me[None] * jac[:, None, None]
            rows.append(np.repeat(conn, nbf, axis=1).ravel())
            cols.append(np.tile(conn, (1, nbf)).ravel())
            vals.append(Mee.ravel())
        M = sp.coo_matrix((np.concatenate(vals),
                           (np.concatenate(rows), np.concatenate(cols))),
                          shape=(dm.n_nodes,) * 2).tocsr()
        return (self.T.T @ M @ self.T).tocsr()

    def load(self, t):
        F_full = wp.zeros(self.dm.n_nodes, dtype=wp.float64, device=DEVICE)
        for pv, b in self.dm.bins.items():
            fq = wp.array(f_ex(self.xq[pv], t), dtype=wp.float64,
                          device=DEVICE)
            lk = make_load_kernel(b["nbf"], b["nqp"], self.dm.dim)
            wp.launch(lk, dim=len(b["eids"]),
                      inputs=[b["conn"], b["h"], b["N"], b["w"], fq, F_full],
                      device=DEVICE)
        return np.asarray(self.T.T @ F_full.numpy())

    def march(self, dt, T_final, order):
        hist = History()
        hist.rotate(u_ex(self.coords, 0.0))          # exact initial data
        t = 0.0
        lus = {}                                     # factorization cache
        while t < T_final - 1e-12:
            t_new = t + dt
            o = bdf_order_now(t_new, dt, order, have_history=hist.have(2))
            b0, b1, b2 = bdf_coeffs(o, dt)
            rhs = self.load(t_new) - self.M @ hist.rhs_history(b1, b2, dt)
            if o not in lus:                         # J changes only with o
                A = ((b0 / dt) * self.M + self.K).tolil()
                for i in self.bdry:
                    A.rows[i] = [int(i)]
                    A.data[i] = [1.0]
                lus[o] = (splu(A.tocsr().tocsc()), None)
            rhs[self.bdry] = u_ex(self.coords[self.bdry], t_new)
            u = lus[o][0].solve(rhs)
            hist.rotate(u, dt=dt)
            t = t_new
        return hist.pre1


if __name__ == "__main__":
    hs = HeatSolver(5)
    T_final = 0.5
    ref = hs.march(T_final / 512, T_final, order=2)   # same-mesh reference
    for order in (1, 2):
        errs = []
        for n in (32, 64, 128):
            u = hs.march(T_final / n, T_final, order)
            errs.append(np.sqrt(((u - ref) ** 2).mean()))
        rates = [np.log2(errs[i] / errs[i + 1]) for i in range(2)]
        print(f"BDF{order}: errors " + "  ".join(f"{e:.3e}" for e in errs)
              + "   rates " + "  ".join(f"{r:.2f}" for r in rates))
\end{lstlisting}

\begin{explore}
\begin{verbatim}
(a) Compare against the EXACT solution instead of the same-mesh
      reference. The BDF2 ladder flattens — at what dt does spatial error
      take over, and does that dt match the h^2 estimate from A1?
  (b) Kill the bootstrap: force BDF2 from step 1 by seeding the second
      history slot with u_ex(t = -dt) (exact history). Does the rate
      change? Now seed it with u_ex(t=0) (WRONG history, O(dt) error).
      Measure the damage.
  (c) Variable steps: alternate dt and 2dt using bdf_coeffs(2, dt, dt_prev).
      Confirm order 2 survives (the production variable-step table).
  (d) PERFORMANCE CORNER: this loop refactorizes NOTHING per step (the
      matrix is constant — we cached the LU). Time a variant that
      refactorizes every step. At level 6, what fraction of the step is
      the factorization? This is why production freezes matrices whenever
      physics allows (and why track-D steppers, whose matrices change
      every step, need iterative solvers).
\end{verbatim}
\end{explore}
